%%% ==========================================================================
%%% Post-mortem: 6-hour Colab Run, 2026-04-14
%%% AI 570 - Team 4 - Wafer Defect Detection
%%% ==========================================================================
\documentclass[conference]{IEEEtran}

\usepackage{amsmath,amssymb}
\usepackage{booktabs}
\usepackage[hidelinks]{hyperref}
\usepackage{float}
\usepackage{url}
\usepackage{listings}
\usepackage{xcolor}

\lstset{
  basicstyle=\ttfamily\footnotesize,
  breaklines=true,
  columns=fullflexible,
  keepspaces=true,
  showstringspaces=false,
  frame=single,
  framesep=4pt,
  rulecolor=\color{gray!40}
}

\begin{document}

\title{Post-mortem: Six-hour Colab Run, 2026-04-14\\
\large{One Headline Result, Two Collapsed Baselines, and a Metrics-File Race}}

\author{
\IEEEauthorblockN{AI 570 -- Team 4}
\IEEEauthorblockA{Anindita Paul, Brandon Pyle, Anand Rajan, Brett Rettura\\
The Pennsylvania State University}
}

\maketitle

\begin{abstract}
A six-hour training run on Google Colab Pro (2026-04-14) trained three
WM-811K wafer defect classifiers in sequence: a custom CNN, ResNet-18, and
EfficientNet-B0. The custom CNN delivered a strong baseline (accuracy 0.961,
macro F1 0.799, ECE 0.007). The two pretrained models---trained under the
same default loss configuration---collapsed: ResNet-18 partially (macro F1
0.659 driven by single-class Edge-Loc confusion), EfficientNet-B0 fully
(macro F1 0.100, validation loss oscillating 20--123 with no convergence).
The aggregated \texttt{metrics.json} artifact captured by Colab and synced
to Google Drive contained metrics for only the last model trained
(EfficientNet), masking the CNN's success behind a metrics-file
race condition. We document both the orchestration bug and the loss-function
configuration that explains why the same training recipe worked for the
custom CNN and failed for the pretrained networks. Both root causes are
remediated in commit \texttt{83ba643}.
\end{abstract}

\begin{IEEEkeywords}
post-mortem, root-cause analysis, class imbalance, deferred re-weighting,
WM-811K, training stability, custom CNN, transfer learning
\end{IEEEkeywords}

\section{Introduction}
The team's PyTorch training pipeline was redesigned the prior day to
support multi-model runs from Google Colab via a subprocess-per-model loop
in the quickstart notebook. The 2026-04-14 run was the first end-to-end
exercise of that loop with cached preprocessing and a Colab Pro T4 runtime.

Wall-clock time totaled approximately six hours: 1.8~h for the custom CNN,
1.8~h for ResNet-18, 1.7~h for EfficientNet-B0. All three model checkpoints
were saved to disk and synced to Google Drive. The aggregated metrics file,
however, contained results for only the last model trained, hiding both the
strong CNN baseline and the diagnostic information that would have flagged
the two pretrained-model collapses.

\section{Headline Result: Custom CNN}

The custom CNN trained without incident and delivered strong test-set
metrics on the standard $70/15/15$ stratified split:

\begin{table}[H]
\centering
\begin{tabular}{lr}
\toprule
Metric & Value \\
\midrule
Accuracy        & 0.9611 \\
Macro F1        & 0.7988 \\
Weighted F1     & 0.9632 \\
Expected Calibration Error & 0.0068 \\
Negative Log Likelihood    & 0.1228 \\
Brier Score                & 0.0578 \\
Wall-clock                 & 6543 s ($\sim$1.8 h) \\
Best epoch (of 10)         & 10 (still improving) \\
\bottomrule
\end{tabular}
\caption{Custom CNN test-set metrics, Colab Pro T4, seed 42.}
\end{table}

Per-class behavior is balanced across rare and common classes:

\begin{table}[H]
\centering
\begin{tabular}{lrrr}
\toprule
Class & Precision & Recall & F1 \\
\midrule
Center      & 0.86 & 0.89 & 0.88 \\
Donut       & 0.58 & 0.87 & 0.70 \\
Edge-Loc    & 0.71 & 0.82 & 0.76 \\
Edge-Ring   & 0.99 & 0.93 & 0.96 \\
Loc         & 0.62 & 0.71 & 0.66 \\
Near-full   & 0.76 & \textbf{1.00} & 0.86 \\
Random      & 0.66 & 0.85 & 0.75 \\
Scratch     & 0.53 & 0.81 & 0.64 \\
none        & 0.99 & 0.98 & 0.99 \\
\bottomrule
\end{tabular}
\caption{Custom CNN per-class metrics. Near-full (22 test samples)
reaches 100\% recall with reasonable precision; Scratch (179 samples)
hits 81\% recall.}
\end{table}

Validation loss declined monotonically from 1.35 to 0.38 over ten epochs;
validation macro F1 climbed from 0.35 to 0.82. The model was still
improving at epoch 10. Training would benefit from a longer schedule.

These numbers compare favorably to published WM-811K baselines and are
adopted as the team's headline result.

\section{ResNet-18: Partial Collapse}

ResNet-18 reached test accuracy 0.777 and macro F1 0.659. Headline-level
metrics look mediocre rather than catastrophic, but per-class confusion
shows a localized failure:

\begin{table}[H]
\centering
\begin{tabular}{lrrr}
\toprule
Class & Precision & Recall & F1 \\
\midrule
\textbf{Edge-Loc} & \textbf{0.14} & \textbf{0.91} & 0.24 \\
Loc               & 0.29 & 0.57 & 0.39 \\
Near-full         & 0.42 & 1.00 & 0.59 \\
Edge-Ring         & 0.92 & 0.96 & 0.94 \\
Center            & 0.80 & 0.91 & 0.85 \\
Random            & 0.80 & 0.86 & 0.83 \\
Scratch           & 0.43 & 0.55 & 0.48 \\
Donut             & 0.87 & 0.65 & 0.74 \\
none              & 0.99 & 0.76 & 0.86 \\
\bottomrule
\end{tabular}
\caption{ResNet-18 per-class metrics. Edge-Loc shows the canonical
over-prediction signature (high recall, very low precision).}
\end{table}

The model has learned to over-predict ``Edge-Loc'' for many true
\texttt{none} samples (driving Edge-Loc recall to 0.91 but precision to
0.14, and \texttt{none} recall down to 0.76). This single-class
mis-allocation accounts for most of the macro F1 deficit. Validation loss
fluctuated between 0.6 and 1.1 across epochs without diverging, so the
collapse is partial rather than catastrophic. The likely fix is the same
as for EfficientNet (Section IV) plus a longer warm-up.

\section{EfficientNet-B0: Full Collapse}

EfficientNet-B0 produced pathological metrics: accuracy 0.374, macro F1
0.100, ECE 0.454. Validation loss oscillated wildly across epochs:

\begin{lstlisting}
val_loss: 29.5, 123.4, 20.5, 80.8, 43.6,
          57.2, 67.0, 78.2, 28.5, 58.4
val_acc:   0.25, 0.37, 0.22, 0.38, 0.25,
           0.26, 0.33, 0.26, 0.25, 0.31
\end{lstlisting}

while training loss declined smoothly from 1.27 to 0.65 and training
accuracy reached 84.6\%. The 100$\times$ separation between train and
validation loss magnitudes is not a typical overfitting signature; it is
a pipeline-level pathology.

Per-class collapse is severe:

\begin{table}[H]
\centering
\begin{tabular}{lrrr}
\toprule
Class & Precision & Recall & F1 \\
\midrule
\textbf{Near-full} & \textbf{0.003} & \textbf{0.864} & 0.005 \\
Random             & 0.025 & 0.585 & 0.048 \\
Scratch            & 0.006 & 0.006 & 0.006 \\
none               & 0.883 & 0.420 & 0.569 \\
\bottomrule
\end{tabular}
\caption{EfficientNet-B0 selected per-class metrics. Near-full reaches
86\% recall with 0.3\% precision: the model predicts ``Near-full'' for
nearly every input.}
\end{table}

\section{Why CNN Worked but Pretrained Models Collapsed}

All three runs used the same loss configuration, the same dataset split,
the same seed, and the same Colab T4. The difference is the learning rate:
the custom CNN trained at $\eta = 10^{-3}$, the two pretrained models at
$\eta = 10^{-4}$ (the standard fine-tuning rate). The class-weight tensor
the loss received was identical in all three runs:

\begin{lstlisting}
Class weights (logged at run time):
  Center=4.47, Donut=34.58, Edge-Loc=3.70,
  Edge-Ring=1.99, Loc=5.35, Near-full=129.34,
  Random=22.20, Scratch=16.11, none=0.13
\end{lstlisting}

That is a 1000$\times$ ratio between Near-full ($w = 129$) and
\texttt{none} ($w = 0.13$). Under inverse-frequency cross-entropy, a
single mispredicted Near-full sample contributes hundreds of units to the
batch-mean loss, while a wrong \texttt{none} sample contributes
$\sim 0.0001$.

The custom CNN's higher learning rate moved its parameters fast enough in
early epochs to find a basin where common-class predictions were already
roughly correct before the rare-class weight gradients dominated. The
pretrained models, which start from already-useful ImageNet features and
therefore need small careful steps, were dragged off the good initial
manifold by the heavy rare-class gradients before they could establish a
common-class baseline. ResNet, with somewhat more capacity, partially
recovered (Edge-Loc collapse but everything else within reason).
EfficientNet, with stronger pretrained inductive biases and a smaller
update step, never escaped the rare-class-driven over-prediction regime.

This is exactly the failure mode Cao et al.~\cite{cao2019} introduced
Deferred Re-Weighting (DRW) to prevent: leave the loss unweighted for the
first $N$ epochs (representation-learning phase), then enable inverse-
frequency weights for the remaining epochs (classifier-adaptation phase).
The trainer in \texttt{src/training/trainer.py} already supports
\texttt{drw\_epoch}, but the configuration default at run time was 0
(disabled).

\textbf{Fix.} \texttt{config.yaml} now ships with
\texttt{training.loss.weighted: false} as the default. The accompanying
comment explains the trade-off and points to a follow-on rare-class study
which benchmarks weighted CE+DRW, focal loss, class-balanced sampling,
and synthetic augmentation against the unweighted baseline on three random
seeds each.

\section{The Metrics-File Race}

Even with one strong baseline and two informative failures, the run's
\texttt{metrics.json} captured by Colab contained only the EfficientNet
entry. The notebook's \texttt{Cell 6} invokes \texttt{train.py} once per
model:

\begin{lstlisting}[language=Python]
for _model in ['cnn', 'resnet', 'efficientnet']:
    subprocess.call(['python', 'train.py',
                     '--model', _model, ...])
\end{lstlisting}

\texttt{train.py} (pre-\texttt{83ba643}) ended each invocation with:

\begin{lstlisting}[language=Python]
metrics_path = self.results_dir / 'metrics.json'
with open(metrics_path, 'w') as f:
    json.dump(self.results, f, indent=2)
\end{lstlisting}

Each subprocess writes \texttt{metrics.json} containing only the model
trained by that subprocess (\texttt{self.results} is initialized empty
per process). The next subprocess overwrites it. Only the last
subprocess's contents survive on disk.

\textbf{Fix.} \texttt{train.py:937} now loads any existing
\texttt{metrics.json}, merges the current process's results, and writes
back. Subprocess loops accumulate correctly. Backward-compatible.

The CNN and ResNet metrics for the 2026-04-14 run were recovered from the
notebook stdout that was preserved in Colab's cell history (the user
copy-pasted them verbatim). The recovered values are committed in
\texttt{results/metrics.json}, with each entry annotated with its source
run and a notes field flagging the collapse cases.

\section{Disagreement with a Parallel Review}

A second agent independently reviewed the same run from the
\texttt{metrics.json} that contained only EfficientNet. It correctly
identified the symptoms (rare-class collapse, validation-loss explosion,
miscalibration) but diagnosed the cause as ``class imbalance not
addressed: \texttt{adaptive\_rebalance: false}, no DRW schedule
(\texttt{drw\_epoch: null})'' and recommended adding class weighting.

This reading is incorrect. The fields visible in
\texttt{epoch\_history}---\texttt{adaptive\_rebalance: false},
\texttt{drw\_epoch: null}---report only the trainer-internal scheduling
flags. They do not indicate whether the underlying CrossEntropyLoss has
class weights attached. As shown in \texttt{train.py:528}, the
auto-computed inverse-frequency \texttt{loss\_weights} tensor is passed to
the criterion whenever \texttt{loss\_cfg.weighted} is true, regardless of
the DRW or adaptive-rebalance flags. The run was using inverse-frequency
weighting; that is precisely the cause of the failure. Adding more
weighting would deepen the collapse.

The peer-review agent's other observations (LR too high for fine-tuning,
EMA disabled, best epoch 4 of 10) are valid follow-ups noted below.

\section{Performance Improvements Landed in the Same Batch}

Three additional changes addressing throughput rather than correctness
landed in commit \texttt{83ba643}:

\begin{itemize}
\item \texttt{training.mixed\_precision: true} (was \texttt{false}). AMP
gives 15--25\% training speedup on CUDA with negligible accuracy cost.
\item \texttt{persistent\_workers=True}, \texttt{prefetch\_factor=2} added
to the DataLoader \texttt{kwargs}. Cuts per-epoch worker respawn overhead
on Colab Pro by 10--20\%.
\item \texttt{set\_seed()} now enables \texttt{cudnn.benchmark=True} by
default with a new \texttt{deterministic=False} kwarg. Yields 8--15\%
speedup on Ampere-class GPUs (Colab Pro A100). Determinism is opt-in for
CI metric-gate runs.
\end{itemize}

Combined: $\sim$30--50\% faster training on Colab Pro. The next 10-epoch
CNN run should complete in approximately one hour instead of 1.8 hours.

\section{Follow-up Items}

\begin{enumerate}
\item Re-train ResNet-18 and EfficientNet-B0 under the new defaults.
Expected behavior: both should converge to macro F1 in the 0.7--0.8 range
matching the custom CNN, possibly higher with proper learning-rate warm-up.
\item Pretrained-model learning rate is still configured at $10^{-4}$
constant. $10^{-5}$ with \texttt{warmup\_epochs > 0} and cosine decay
is more standard for fine-tuning.
\item Enable EMA (\texttt{use\_ema: true}) for pretrained-model runs to
smooth noisy validation metrics.
\item Run the rare-class study (\texttt{docs/rare\_class\_study.md}) to
benchmark each rebalancing intervention head-to-head. The custom CNN's
strong rare-class recall (Near-full 1.00, Scratch 0.81, Donut 0.87)
without DRW or balanced sampling is itself an interesting baseline.
\end{enumerate}

\section{Conclusion}

The run produced one strong, reportable baseline (custom CNN, accuracy
0.961, macro F1 0.799) and two informative failure modes (ResNet partial
collapse, EfficientNet full collapse). The aggregated metrics artifact
captured only the worst result, masking the success behind a
metrics-file race condition. Both the orchestration bug and the
configuration default that drove the pretrained-model collapses are
remediated in commit \texttt{83ba643}.

The custom CNN result stands as the team's headline number. The two
pretrained-model checkpoints are not usable as reported baselines and
should be re-trained under the new defaults.

\begin{thebibliography}{1}

\bibitem{cao2019}
Y. Cao, C. Wei, A. Gaidon, N. Arechiga, and T. Ma,
``Learning Imbalanced Datasets with Label-Distribution-Aware Margin Loss,''
in \emph{Advances in Neural Information Processing Systems}, 2019.
arXiv:1906.07413.

\bibitem{wuwm811k}
M.-J. Wu, J.-S. R. Jang, and J.-L. Chen,
``Wafer Map Failure Pattern Recognition and Similarity Ranking for
Large-Scale Data Sets,''
\emph{IEEE Transactions on Semiconductor Manufacturing},
vol. 28, no. 1, pp. 1--12, 2015.

\end{thebibliography}

\end{document}
